\section{Modeling Landscape and Design Dimensions}

Detection and preictal prediction from ECG/HRV, PPG, EDA, and ACC are
framed by windowing, labeling, and alarm logic, with alarm-based
sensitivity, FAR/h, and latency as the outcomes used in the research
questions.

\subsection{Architectural Evolution: Feature Engineering vs. Deep Learning}
Classical pipelines utilizing HRV, EDA, and ACC features combined with
SVM, RF, or kNN remain competitive in retrospective settings
\parencite{sethFeasibilityCardiacbasedSeizure2023,masonHeartRateVariability2024}.
However, deep temporal models such as 1D-CNN, LSTM, and TCN often
improve phase discrimination. Recently, attention variants and
Transformers are being explored on these signals
\parencite{yangMultimodalAISystem2022,pordoyEnhancedNonEEGMultimodal2025,
abtahiIdentificationRelevantECG2025,kalousiosECGbasedEpilepticSeizure2024}.
A key inquiry of this review is determining how feature-based models
compare to deep models and whether attention mechanisms add distinct
value under ambulatory, alarm-based evaluation.

\subsection{Sensor Modalities and Fusion Strategies}
Fusion across ECG, EDA, ACC, and PPG at the feature, decision, or
representation level seeks complementary cues and is relevant to
maintaining performance under noise and non-stationarity
\parencite{yangMultimodalAISystem2022,pordoyEnhancedNonEEGMultimodal2025}.

\subsection{Personalization and Domain Shift}
Patient-specific fine-tuning and post-hoc calibration
shift the operating point toward lower FAR/h and higher sensitivity in
ambulatory use \parencite{andradePerformanceSeizurePrediction2024}.
Thresholding and periodic recalibration are commonly used to handle
drift.

\subsection{Robustness and Generalization}
Robustness requires explicit checks of domain shift across patient,
device, site, and activity. Cross-dataset evaluation (e.g., training
on TUH and testing on RPAH or EPILEPSIAE) exposes this shift and favors
alarm-based over sample-based metrics
\parencite{yangMultimodalAISystem2022,andradePerformanceSeizurePrediction2024}.
Artifact handling, augmentation through time-warping, jitter, or
scaling, and Out-of-Distribution (OOD) detection support
generalization \parencite{kalousiosECGbasedEpilepticSeizure2024,
sethFeasibilityCardiacbasedSeizure2023}.

\subsection{Deployment Constraints: Latency, Energy, and Interpretability}
Wearable seizure detection imposes strict deployment constraints.
Detection latency must remain below clinically relevant thresholds 
(typically seconds to allow caregiver response or patient warning), 
which limits model complexity and inference time.
Energy consumption determines battery life; 
continuous inference on resource-constrained MCUs or mobile processors 
requires lightweight architectures such as quantized CNNs or efficient RNNs.
Compact sequence models (e.g., depthwise-separable 1D-CNNs, temporal convolutional networks) 
and event-driven or on-demand processing reduce computational burden.

Interpretability remains critical for clinical adoption. 
Regulators and clinicians require transparency in alarm generation 
to trust automated decisions affecting patient safety.
Post-hoc explanation methods (e.g., attention weights, Grad-CAM on temporal segments, 
SHAP values for HRV features) can highlight which signal segments or features triggered an alarm.
Hybrid pipelines that combine interpretable HRV/EDA features with deep representations 
offer a middle ground between performance and explainability.
